\label{sec-intro}
\section{Introduction}
Fake news has become a medium by which to spread
misinformation~\cite{oshikawa-etal-2020-survey, acmwebfakenews}.
One common way to fight against fake news is to release verified
news.\footnote{In this work, we define ``verified news'' as news written
specifically to clarify false information; the term ``real news''  is defined
as general news that does not contain misinformation.} However,
as the goal of news verification is to correct misinformation, verified news headlines
are often bland, making it difficult to gain the attention of users, which
works against the need to alleviate the harmful impact of fake news.
Therefore, headlines for verified news articles should be rewritten to be more
intriguing but still faithful, which is expected to pique reader interest in
verified news. Many studies have been
conducted on generating attractive
headlines~\cite{jin-etal-2020-hooks,xu-etal-2019-clickbait}, among which
\textit{clickbait} represents the style that generates the most reads or
clicks.
Despite their success in attracting readers, there are several challenges in
current models. 
First, clickbait datasets for training headline generators with sensational
style transfer are commonly collected based on the amount of views or clicks,
which assumes that headline popularity is always due to 
the writing style~\cite{popularitybased}. However, user reading preferences could also be motivated by
trending topics or major events. For instance, ``\textit{Flights cancelled as
typhoon nears}'' was the most popular news on a day that a typhoon was coming.
Although such headlines get many views and clicks, the writing style itself is
not interesting, and could end up as noise
in the dataset. 
Second, harmful ``hallucinations'' created by headlines exaggerated to be more
sensational could distort the meaning of the original article. This is
especially critical as we do not want our model itself to spread
misinformation. 
However, as such sensational headline generation models often generate
clickbait with more ambiguous words, it increases the difficulty of evaluating faithfulness by aligning
title semantics with the news content. 

In this work, we propose making real news intriguing by learning what fake
news is good at. We seek to learn what makes fake news eye-catching
instead of simply mimicking the titles of fake news. 
Quantity-wise, the many circulating fake news articles serve as learning materials by
which we can learn to generate more attractive headlines; 
style-wise, fake news is deliberately written to attract attention.  
To learn such attractive writing styles, we adopt the forward-reference (FR)
writing technique~\cite{BLOM201587}, which draws from
psychology and journalism, and is frequently used to create attractive headlines.
Specifically, FR creates an information gap between readers and the news content
with the headline, motivating the reader's curiosity~\cite{psychologyofcuriosity} 
to investigate the news content, and hence provoking the desire to click on the
headline.
One example is the headline \textit{``Wanna be an enviable couple? 12 things a
happy couple must do... It's that simple!''}, which drives readers to find out
what those things are. 

Here, to understand the relation between veracity, attractiveness, and FR types
in news headlines, we conducted a preliminary user study to investigate the
attractiveness of fake and real news, and analyzed the FR types
used in headlines in terms of veracity.
Given these results and observations, we propose \modelname, a novel framework
by which to generate attractive but faithful headlines. 
In this framework, we use FR to remove the need to learn directly from the
click-based dataset.
%As we mentoined above, evaluating faithfulness of clickbait headlines is challenging in many ways~\cite{onfaith, kryscinski-etal-2020-evaluating}
%, as such headlines usually contains more fuzzy information. 
%So a soft constraint to evaluate semantic alignment is more suitable than hard raw metrics like ROUGE~\cite{lin-2004-rouge}. 
To ensure the faithfulness of the generated headlines, we design a
lexical-bias-robust textual entailment component on the generated headline and
its original content to confirm that the content infers the headline. 
%We then use this as a discriminator to calculate the faithfulness reward for each training step. 
In addition, we propose \dataset, an innovative dataset which consists of pairs
of fake and verified news headlines, their content, and their FR types. We
conduct experiments on \dataset and evaluate the results in terms of both automatic
metrics and human evaluation. 
%Results show superiority of our method to generate more faithful headline, while the curiosity motive is competitive to method that is designed specifically for sensational headline generation.
%is effective for generating both faithful and curiosity motive headlines.  
%Finally, we report our ablation study for each component of \modelname, indicating that our two pre-trained scorers have significant improvement on both sides.  
In sum, the contributions of our work are threefold:
\begin{itemize}
	 \item We conduct a thorough user study to understand the relation between
	 reading preferences and FR types on fake news and verified news. 
	 \item We propose a novel framework for generating attractive but faithful
	 headlines. In human evaluations, \modelname 
	  largely outperforms baselines on attractiveness and faithfulness.
	 \item We propose a new dataset containing pairs of fake and 
	 verified news, including their headlines, content, and FR
	 types in headlines.
\end{itemize}

% \begin{figure}[t]
% \caption{An illustration of our goal: making title more fuzzy to raise users' curiosity. We first extract entities from the original title, and then extend the attribute graph by querying ConceptNet with the extracted entities. Our model learns to select proper entities to replace the original ones, while it can generate information gap to attract readers. The example attribute graph is correspond to \textit{coffee}, and dots with the same color indicate that they have same relations with \textit{coffee}.} 
% \includegraphics[width=0.48\textwidth]{figures/fig_intro.pdf}
% \centering
% \end{figure}
% % summarize and contribution
